\documentclass[12pt]{article}
\usepackage[T1]{fontenc}
\usepackage[utf8]{inputenc}
\usepackage{lmodern}
\usepackage{textcomp}
\usepackage{enumitem}
\usepackage{geometry}
\usepackage{graphicx}
\usepackage{amsmath, amssymb}
\geometry{margin=1in}
\begin{document}
\begin{center}
Economics - Undergraduate Level Assessment 2 \
Total Marks: 40 \
Answer all questions.
\end{center}

\section*{Multiple Choice (10 marks)}
\begin{enumerate}[label=\arabic*.]
\item If Matt's total utility from consuming bratwurst increased at a constant rate, no matter how many bratwurst Matt consumed, what would Matt's demand curve for bratwurst look like?
\begin{enumerate}[label=(\alph*)]
    \item Linearly increasing
    \item Inverted U-shape
    \item Horizontal
    \item Upward sloping
    \item Vertical
\end{enumerate}
\item Enumerate the cause of the long-and short-run farm problems.
\begin{enumerate}[label=(\alph*)]
    \item High elasticity of demand and supply
    \item Decline and instability
    \item Lack of technological progress
\end{enumerate}
\item What does the Harrod-Domar model, taken with the Keynesian theory of savings, imply for the growth rates of poor countries?
\begin{enumerate}[label=(\alph*)]
    \item Rich and poor countries will experience the same growth rates regardless of savings rates
    \item Poor countries will have fluctuating growth rates independent of savings
    \item The level of technology determines growth rates, not savings
    \item The level of income has no impact on savings
    \item Savings have no correlation with growth rates in any country
\end{enumerate}
\item The real interest rate is
\begin{enumerate}[label=(\alph*)]
    \item the anticipated inflation rate plus the nominal interest rate
    \item the nominal interest rate multiplied by the inflation rate
    \item what one sees when looking at bank literature
    \item the anticipated inflation rate minus the nominal interest rate
    \item the nominal interest rate minus anticipated inflation
\end{enumerate}
\item An increase in corporate optimism will have which of the following effects in the market for loanable funds?
\begin{enumerate}[label=(\alph*)]
    \item A decrease in both supply and demand and an ambiguous change in interest rates.
    \item An increase in supply lowering the interest rate.
    \item A decrease in demand increasing the interest rate.
    \item An increase in demand increasing the interest rate.
    \item A decrease in supply increasing the interest rate.
\end{enumerate}
\end{enumerate}

\section*{True / False (10 marks)}
\textit{Answer True or False}
\begin{enumerate}[label=\arabic*.]
\setcounter{enumi}{5}
\item True or False: Stagflation can occur as a result of a leftward shift in the aggregate supply curve, leading to higher prices and lower output.
\item True or False: An increase in real GDP per capita of \$1,000 over ten years indicates a standard of living increase of 2.5 percent.
\item True or False: In the long run, perfect competition leads to price being greater than marginal cost while equating to average total cost (P \textgreater{} MC = ATC).
\item True or False: For an autoregressive process to be stationary, the roots of the characteristic equation must all be greater than one in absolute value.
\item True or False: Negative externalities arise only when total welfare is maximized in a market.
\end{enumerate}

\section*{Long Form Response (20 marks)}
\textit{Answer each question with a clear and complete explanation.}
\begin{enumerate}[label=\arabic*.]
\setcounter{enumi}{10}
\item Discuss the significance of the relationship between price and marginal cost in monopolistic competition, and analyze how it impacts market outcomes and firm behavior.
\item Discuss how a poll tax of \$2 per head exemplifies a regressive tax system, considering its impact on individuals with varying income levels and the implications for economic equity.
\end{enumerate}

\vfill
\noindent\textit{End of Paper}
\end{document}